\clearpage
\begingroup
\fancyhead[L]{\small POSTHUMOUS PAPERS AND CORRESPONDENCE}
\renewcommand{\headrulewidth}{0pt}
\par\medskip
\SourceLeaf{589}
\Needspace{8\baselineskip}
\pdfbookmark[2]{[4] Nicolai to Gauss, 22 April 1814}{pallas-VIII-letter4}
\begin{center}
[4.]\par\medskip
\textsc{Nicolai} to \textsc{Gauss}, Seeberg, 22 April 1814.
\end{center}
\par\noindent\hspace*{\parindent}\{.\ .\ .\ .\ .\quad You will have received my last letter of 8 April with the perturbation formulas for
the eccentricity and perihelion. Today I can send you
the rest of the work, namely the equations for the 2\textsuperscript{nd}
part of the epoch.\quad .\ .\ .\ .\ .

\par\noindent\hspace*{\parindent}In the analytic expression for \(\mathrm{d}L\) I used a small transformation
that I would like to send you for your consideration.

\par\noindent\hspace*{\parindent}If we set \(1+\cos v\cos E=u\), we obtain
\[
\begin{aligned}
\mathrm{d}\varphi&=-maa\sin v\cos E\mathord{.}\,X+maauY\\
\mathrm{d}\tilde\omega&=A-maa\operatorname{cotang}\varphi\mathord{.}\,(3-u)\mathord{.}\,X
  +\frac{maa}{e}\cos v\sin E\mathord{.}\,Y.
\end{aligned}
\]

\par\noindent\hspace*{\parindent}If the latter, with \(A\) left out, is multiplied by \(1-\dfrac{rr}{aa\cos\varphi}\), and the former by
\(-\left(\cos\varphi+\dfrac{r}{a\cos\varphi}\right)\sin E\), it follows immediately that
\[
\begin{aligned}
\mathrm{d}L={}&A+\frac{m}{e}\bigl\{(rr-aa\cos\varphi)(3-u)+ae(p+r)(2-u)\cos E\bigr\}X\\
 &{}-\frac{m\sin E}{e\cos\varphi}\bigl\{(rr-aa\cos\varphi)\cos v+ae(p+r)u\bigr\}Y.
\end{aligned}
\]

\par\noindent\hspace*{\parindent}I have not tried to simplify the factors in braces any further,
because they seemed convenient enough for numerical calculation, especially with
your auxiliary table. This transformation and the introduction of the
quantity \(u\), incidentally, are based on a property I recently came across by chance
in the relation between the true and eccentric anomalies.
If we regard the complements of \(v\) and \(E\) to \(90^0\)
(or to \(450^0\)) as the two sides of a spherical triangle that enclose the
eccentricity angle \(\varphi\), then the third side of the triangle, opposite
\(\varphi\), is always constant and is itself \(=\varphi\). The two
other angles of this triangle also have curious properties.
The one opposite the side \(\genfrac{}{}{0pt}{0}{90^0}{450^0}-v\) is always equal to that side, and
the other, opposite the side \(\genfrac{}{}{0pt}{0}{90^0}{450^0}-E\), is the complement of that
side to \(180^0\), or \(=90^0+E\).\quad .\ .\ .\ .\ .

\SourceLeaf{590}
\par\noindent\hspace*{\parindent}If our factors for the perturbing forces should fail to agree exactly,
this will be because you have probably used the latest
mean elements, whereas I used those obtained in
the very first calculation by maniples.
\par\medskip
\Needspace{13\baselineskip}
\begin{center}
Second part of the epoch.\par
Annual change \(=-1\overset{\prime\prime}{,}31973\)
\end{center}
\begingroup
\small
\newcommand{\VIIIpallas}{\text{\AstroSymbols ⚴}}
\newcommand{\VIIIsaturn}{\text{\AstroSymbols ♄}}
\setlength{\arraycolsep}{2pt}
\renewcommand{\arraystretch}{1.25}
\[
\begin{array}{r@{\,}l@{\qquad}r@{\,}l}
0\overset{\prime\prime}{,}45&\sin(\VIIIpallas+180^0\ 35')&+0\overset{\prime\prime}{,}27&\sin(3\VIIIpallas-2\VIIIsaturn+65^0\ 10')\\
+0{,}29&\sin(2\VIIIpallas+181\ 47)&+0{,}30&\sin(3\VIIIsaturn+207\ 41)\\
+1{,}78&\sin(\VIIIsaturn+140\ 25)&+1{,}05&\sin(\VIIIpallas-3\VIIIsaturn+163\ 1)\\
+0{,}32&\sin(\VIIIpallas+\VIIIsaturn+181\ 12)&+0{,}45&\sin(2\VIIIpallas-3\VIIIsaturn+154\ 50)\\
+0{,}41&\sin(\VIIIpallas-\VIIIsaturn+178\ 5)&+0{,}42&\sin(3\VIIIpallas-3\VIIIsaturn+173\ 25)\\
+0{,}11&\sin(2\VIIIpallas-\VIIIsaturn+158\ 28)&+0{,}21&\sin(\VIIIpallas-4\VIIIsaturn+166\ 29)\\
+2{,}01&\sin(2\VIIIsaturn+221\ 24)&+0{,}15&\sin(2\VIIIpallas-4\VIIIsaturn+159\ 55)\\
+3{,}35&\sin(\VIIIpallas-2\VIIIsaturn+124\ 44)&+0{,}11&\sin(3\VIIIpallas-4\VIIIsaturn+156\ 39).\}\\
+2{,}04&\sin(2\VIIIpallas-2\VIIIsaturn+178\ 31)&&
\end{array}
\]
\endgroup
\clearpage
\endgroup
